\documentclass[10pt]{article}
\usepackage[utf8]{inputenc}
\usepackage[english]{babel}
 
\usepackage{multicol}
\usepackage[none]{hyphenat}
\usepackage{color}
\usepackage{xcolor}
\usepackage{comment}
\usepackage{tikz}
\usepackage{alltt}
\usetikzlibrary{backgrounds,calc}
\usepackage{minted}
\usepackage[framemethod=TikZ]{mdframed}
\usepackage{wallpaper}
\usepackage{parskip}
\setlength{\parindent}{0pt}
\usepackage[a4paper, margin={1.5cm}]{geometry}
\usepackage{titlesec}
\usepackage{fancyhdr}
\pagestyle{fancy}
\pagenumbering{gobble}


% Github image
\newcommand{\github}{%
\begin{tikzpicture}[y=0.3pt, x=0.3pt,yscale=-1, inner sep=0pt, outer sep=0pt,opacity=0.4]
  \begin{scope}[shift={(506.69823,386.92617)}]
    \path[fill=black] (116.9933,59.7217) .. controls (116.9933,71.2283) and
      (107.6655,80.5562) .. (96.1589,80.5562) .. controls (84.6524,80.5562) and
      (75.3245,71.2283) .. (75.3245,59.7217) .. controls (75.3245,48.2152) and
      (84.6524,38.8873) .. (96.1589,38.8873) .. controls (107.6654,38.8873) and
      (116.9933,48.2152) .. (116.9933,59.7217) -- cycle;
    \path[cm={{0.88462,0.0,0.0,0.88462,(11.09526,6.89097)}},fill=white]
      (116.9933,59.7217) .. controls (116.9933,71.2283) and (107.6655,80.5562) ..
      (96.1589,80.5562) .. controls (84.6524,80.5562) and (75.3245,71.2283) ..
      (75.3245,59.7217) .. controls (75.3245,48.2152) and (84.6524,38.8873) ..
      (96.1589,38.8873) .. controls (107.6654,38.8873) and (116.9933,48.2152) ..
      (116.9933,59.7217) -- cycle;
    \path[fill=black,nonzero rule] (103.4671,45.2878) .. controls (102.9322,45.4374)
      and (101.2003,46.2576) .. (100.5403,46.6739) -- (100.1099,46.9454) --
      (99.4882,46.8019) .. controls (99.0810,46.7080) and (98.1204,46.6415) ..
      (96.7048,46.6094) .. controls (94.4953,46.5593) and (93.4339,46.6361) ..
      (92.2380,46.9324) -- (91.6450,47.0793) -- (90.9468,46.6426) .. controls
      (90.0955,46.1101) and (88.7784,45.4948) .. (88.1825,45.3512) .. controls
      (87.9348,45.2916) and (87.5225,45.2429) .. (87.2643,45.2429) .. controls
      (86.8530,45.2429) and (86.7816,45.2733) .. (86.6817,45.4916) .. controls
      (86.3049,46.3144) and (86.1702,48.1697) .. (86.3982,49.3940) --
      (86.5087,49.9870) -- (86.0485,50.6088) .. controls (85.4184,51.4600) and
      (84.9876,52.3958) .. (84.8509,53.2104) .. controls (84.6439,54.4443) and
      (84.8398,57.3849) .. (85.1880,58.2702) .. controls (85.2564,58.4443) and
      (85.2939,58.4403) .. (81.6976,58.6338) .. controls (79.2203,58.7672) and
      (77.4880,58.9815) .. (77.2948,59.1788) .. controls (77.1683,59.3080) and
      (77.2021,59.3161) .. (77.6325,59.2604) .. controls (79.8802,58.9695) and
      (83.0680,58.7293) .. (84.6818,58.7293) .. controls (85.3322,58.7293) and
      (85.3437,58.7337) .. (85.4709,59.0402) .. controls (85.5424,59.2123) and
      (85.5936,59.3574) .. (85.5857,59.3654) .. controls (85.5778,59.3733) and
      (84.8826,59.4288) .. (84.0409,59.4888) .. controls (82.1375,59.6245) and
      (80.3024,59.8884) .. (78.6942,60.2577) .. controls (77.5177,60.5279) and
      (77.1884,60.6573) .. (77.3264,60.7953) .. controls (77.3578,60.8267) and
      (77.9386,60.7190) .. (78.6081,60.5575) .. controls (80.6932,60.0548) and
      (83.4463,59.6858) .. (85.1122,59.6858) .. controls (85.7817,59.6858) and
      (85.8050,59.6938) .. (85.9497,59.9727) .. controls (86.1509,60.3606) and
      (87.1973,61.4638) .. (87.6756,61.7923) .. controls (88.7575,62.5354) and
      (90.1146,63.0487) .. (91.7311,63.3262) .. controls (92.3241,63.4280) and
      (92.8529,63.5117) .. (92.9028,63.5117) .. controls (92.9519,63.5117) and
      (92.8171,63.7221) .. (92.6084,63.9708) .. controls (92.2151,64.4395) and
      (91.8427,65.1574) .. (91.8393,65.4534) .. controls (91.8343,65.8877) and
      (90.1911,66.2247) .. (89.1390,66.0071) .. controls (88.4365,65.8618) and
      (87.9449,65.5203) .. (87.3370,64.7552) .. controls (86.5997,63.8274) and
      (86.0013,63.2318) .. (85.6000,63.0268) .. controls (85.1313,62.7874) and
      (84.1718,62.7744) .. (83.9782,63.0048) .. controls (83.8657,63.1387) and
      (83.8975,63.1954) .. (84.2322,63.4586) .. controls (85.1908,64.2122) and
      (85.6680,64.7934) .. (86.1681,65.8169) .. controls (86.7336,66.9742) and
      (87.2885,67.5731) .. (88.1825,67.9913) .. controls (88.6992,68.2330) and
      (88.8042,68.2463) .. (90.1911,68.2463) -- (91.6546,68.2463) --
      (91.6259,70.0923) -- (91.5972,71.9383) -- (91.2050,72.2922) .. controls
      (90.9850,72.4908) and (90.6785,72.7603) .. (90.5068,72.9061) .. controls
      (90.0483,73.2955) and (90.1529,73.4104) .. (90.9946,73.4418) .. controls
      (91.6450,73.4662) and (91.7691,73.4390) .. (92.3241,73.1503) .. controls
      (93.3630,72.6098) and (93.3667,72.5983) .. (93.3667,69.8628) .. controls
      (93.3667,67.3377) and (93.4455,66.7059) .. (93.8107,66.3047) --
      (94.0458,66.0464) -- (93.9980,69.2506) .. controls (93.9695,71.1540) and
      (93.9075,72.6024) .. (93.8449,72.8183) .. controls (93.7868,73.0192) and
      (93.6134,73.3252) .. (93.4575,73.5022) .. controls (93.3059,73.6744) and
      (93.1754,73.9155) .. (93.1754,74.0235) .. controls (93.1754,74.1976) and
      (93.2328,74.2243) .. (93.6058,74.2243) .. controls (94.3519,74.2243) and
      (95.3191,73.5586) .. (95.6209,72.8374) .. controls (95.8285,72.3417) and
      (95.9492,70.6280) .. (95.9492,68.1794) -- (95.9492,65.9029) --
      (96.4179,65.9029) -- (96.4465,69.1311) .. controls (96.4752,72.3544) and
      (96.4756,72.3599) .. (96.7144,72.8374) .. controls (97.1209,73.6505) and
      (98.5189,74.4873) .. (99.0195,74.2173) .. controls (99.2785,74.0776) and
      (99.2470,73.9374) .. (98.8154,73.3061) .. controls (98.5996,72.9905) and
      (98.3935,72.5452) .. (98.3372,72.2731) .. controls (98.2088,71.6514) and
      (98.2544,66.1949) .. (98.3882,66.1752) .. controls (98.4417,66.1673) and
      (98.5682,66.3047) .. (98.6752,66.4864) .. controls (98.8508,66.7849) and
      (98.8704,67.0316) .. (98.9143,69.4898) .. controls (98.9477,71.3645) and
      (98.9985,72.2310) .. (99.0833,72.3783) .. controls (99.2883,72.7344) and
      (99.9568,73.2398) .. (100.3777,73.3570) .. controls (100.6002,73.4189) and
      (101.0568,73.4562) .. (101.4011,73.4406) .. controls (102.2046,73.4043) and
      (102.2524,73.2299) .. (101.5924,72.7428) .. controls (100.6531,72.0496) and
      (100.6840,72.1775) .. (100.6746,68.9637) .. controls (100.6656,65.9699) and
      (100.6109,65.4703) .. (100.2007,64.6499) .. controls (100.0812,64.4108) and
      (99.8134,64.0644) .. (99.5982,63.8704) -- (99.2108,63.5213) --
      (99.6603,63.4617) .. controls (100.5690,63.3414) and (102.0372,63.0328) ..
      (102.6446,62.8345) .. controls (104.1654,62.3382) and (105.5084,61.3208) ..
      (106.1445,60.1832) -- (106.4227,59.6858) -- (106.9679,59.6858) .. controls
      (108.9956,59.6858) and (112.7503,60.2177) .. (114.7632,60.7901) .. controls
      (114.9162,60.8337) and (114.9832,60.8090) .. (114.9832,60.7092) .. controls
      (114.9832,60.3420) and (111.4059,59.7105) .. (108.1061,59.4950) .. controls
      (107.2931,59.4419) and (106.6181,59.3838) .. (106.5996,59.3654) .. controls
      (106.5815,59.3473) and (106.6145,59.1932) .. (106.6713,59.0306) --
      (106.7765,58.7293) -- (107.9817,58.7323) .. controls (109.6496,58.7363) and
      (111.7789,58.8872) .. (113.5293,59.1252) .. controls (114.8684,59.3073) and
      (115.2129,59.3130) .. (115.0501,59.1502) .. controls (114.8456,58.9456) and
      (112.1137,58.6482) .. (109.3399,58.5285) .. controls (108.0008,58.4707) and
      (106.8944,58.4168) .. (106.8865,58.4089) .. controls (106.8785,58.4010) and
      (106.9394,58.0694) .. (107.0204,57.6772) .. controls (107.1184,57.2030) and
      (107.1719,56.3764) .. (107.1782,55.2382) .. controls (107.1862,53.7174) and
      (107.1624,53.4295) .. (106.9708,52.7704) .. controls (106.6953,51.8235) and
      (106.3173,51.0734) .. (105.7225,50.2931) -- (105.2557,49.6810) --
      (105.2940,48.0598) .. controls (105.3295,46.5581) and (105.3160,46.3927) ..
      (105.1123,45.8168) -- (104.8923,45.1951) -- (104.4140,45.1760) .. controls
      (104.1462,45.1653) and (103.7296,45.2145) .. (103.4671,45.2879) -- cycle;
  \end{scope}
\end{tikzpicture}
}

\lfoot{
\github https://github.com/anchor/haskell-cheat-sheets
}

% spacing
\titlespacing*{\subsection}{0pt}{3pt}{5pt}

% Anchor colors
\definecolor{anchorstripe}{HTML}{50a0e0}
\definecolor{anchorbase}{HTML}{356895}
\definecolor{anchorextra}{HTML}{27ccc0}
\definecolor{anchoraccent}{HTML}{F78224}
\definecolor{shadecolor}{HTML}{27ccc0}
 
% Custom commands
\newcommand{\HRule}[1][\medskipamount]{\par
  \vspace*{\dimexpr-\parskip-\baselineskip+#1}
  \noindent\rule{\linewidth}{0.1mm}\par
  \vspace*{\dimexpr-\parskip-.5\baselineskip+#1}}

\newmdenv[%
	linecolor=anchorextra,%
	backgroundcolor=anchorextra!1,%
	outerlinewidth = 1,%
	roundcorner=2pt,%
]{box2}

\newmdenv[%
	linecolor=anchoraccent,%
	backgroundcolor=anchoraccent!1,%
	outerlinewidth = 1,%
	roundcorner=2pt,%
]{box1}




% Spacing for logo
\addtolength{\wpXoffset}{5.5cm}
\addtolength{\wpYoffset}{13.1cm}

\begin{document}

\CenterWallPaper{0.3}{common/anchor-logo.png}


{\huge \bfseries Control.Lens (viewing) \\[0.2cm]}

\HRule%

\begin{multicols}{2}

\begin{box1}
\subsection *{Getting with Getters}

Any function $(s \to a)$ can be flipped into continuation passing style, $(a \to r)
\to s \to r$ and decorated with \textbf{Const} to obtain:

\begin{minted}{haskell}
type Getting r s a =
  (a -> Const r a) -> s -> Const r s
\end{minted}
A \textbf{Getter} describes how to retrieve a single value in a way that can be composed
with other \textbf{LensLike} constructions.

When you see this in a type signature it indicates that you can pass the
function a \textbf{Lens}, \textbf{Getter}, \textbf{Traversal}, \textbf{Fold},
\textbf{Prism}, \textbf{Iso}, or one of the indexed
variants, and it will just ``do the right thing''.
\end{box1}

\begin{box2}
\subsection *{Safe head}
Perform a safe head of a \textbf{Fold} or \textbf{Traversal} or retrieve
\textbf{Just} the result from a \textbf{Getter} or \textbf{Lens}.

$ (\ \hat{}\,?) \equiv flip\;preview $

\begin{minted}{haskell}
(^?) :: s -> Getting (First a) s a -> Maybe a

> Right 4 ^?_Left
Nothing
> "world" ^? ix 3
Just `l'
\end{minted}
\end{box2}

\begin{box1}

\subsection *{Viewing lenses}

View the value pointed to by a \textbf{Getter} or \textbf{Lens} or the result
of folding over all the results of a \textbf{Fold} or \textbf{Traversal} that
points at a monoidal values.

This is the same operation as \textbf{view} with the arguments flipped.

\begin{minted}{haskell}
(^.) :: s -> Getting a s a -> a 

> (0, -5)^._2.to abs
5
> ["a", "b", "c"] ^. traversed
"abc"
\end{minted}
\end{box1}

\begin{box2}
\subsection*{Using MonadState}

Use the target of a \textbf{Lens}, \textbf{Iso}, or \textbf{Getter} in the
current state, or use a summary of a \textbf{Fold} or \textbf{Traversal} that
points to a monoidal value.

\begin{minted}{haskell}
use :: MonadState s m => Getting a s a -> m a

> evalState (use _1) (1,2)
1
> evalState (uses _1 length) ("hello","")
5
\end{minted}
\end{box2}

\columnbreak

\begin{box1}
\subsection *{Folding Foldables}
\begin{minted}{haskell}
type Fold s a =
  forall m. Monoid m => Getting m s a
\end{minted}
A \textbf{Fold s a} is a generalization of something \textbf{Foldable}. It
allows you to extract multiple results from a container. Every \textbf{Getter}
is a valid \textbf{Fold} that simply doesn't use the Monoid it is passed.

If there exists a \textbf{foo} method that expects a
\textbf{Foldable (f a)}, then there should be a \textbf{fooOf} method that
takes a \textbf{Fold s a} and a value of type \textbf{s}.

\end{box1}

\begin{box2}
\subsection *{Extracting lists from Folds}

Extract a list of the targets of a \textbf{Fold}, an infix version of
\textbf{toListOf}.

$ toList\;xs \equiv xs\;\hat{}.. folded $

\begin{minted}{haskell}
(^..) :: s -> Getting (Endo [a]) s a -> [a] 

> [[1,2],[3]] ^.. traverse . traverse
[1,2,3]
> (1,2) ^.. both
[1,2]
\end{minted}
\end{box2}

\begin{box1}
\subsection *{Checking for matches}

Check to see if this \textbf{Fold} or \textbf{Traversal} matches 1 or more entries. For the
opposite, use \textbf{hasn't}.

\begin{minted}{haskell}
has :: Getting Any s a -> s -> Bool

> has (element 0) []
False
> has _Right (Left 12)
False
> hasn't _Right (Left 12)
True
\end{minted}
\end{box1}


\begin{box2}
\subsection*{Indexed Getters}
For most operations, there is an indexed variant which will work as expected if
the underlying target supports a notion of \textbf{Indexing}.

\begin{minted}{haskell}
> ["ab", "c"] ^@.. itraversed <.> itraversed  
[((0,0),'a'),((0,1),'b'),((1,0),'c')]
> "hello" ^@.. itraversed . indices even
[(0,'h'),(2,'l'),(4,'o')]

> ifind (\i k -> i > k) [1,2,2,2]
Just (3,2)
\end{minted}
\end{box2}
\end{multicols}
\newpage

{\huge \bfseries Control.Lens (setting) \\[0.2cm]}

\HRule%

\begin{multicols}{2}

\begin{box1}
\subsection *{Modifying records with Setters}
A \textbf{Setter s t a b} is a generalization of fmap from \textbf{Functor}. It
allows you to map into a structure and change out the contents, but it isn't
strong enough to allow you to enumerate those contents. Starting with
$fmap :: Functor f \Rightarrow (a \to b) \to f a \to f b$ we monomorphize the type
to obtain $(a \to b) \to s \to t$ and then decorate it with Identity to obtain:

\begin{minted}{haskell}
type Setter s t a b =
  (a -> Identity b) -> s -> Identity t
\end{minted}

Every \textbf{Traversal} is a valid \textbf{Setter}, since \textbf{Identity} is
\textbf{Applicative}.

\end{box1}

\begin{box2}
\subsection *{Modifying with a function}

Modifies the target of a \textbf{Lens} or all of the targets of a
\textbf{Setter} or \textbf{Traversal} with a user supplied function.

This is an infix version of \textbf{over}.

\begin{minted}{haskell}
(%~) :: Profunctor p
     => Setting p s t a b -> p a b -> s -> t

> traverse %~ even $ [1,2,3]
[False,True,False]
\end{minted}
\end{box2}

\begin{box1}
\subsection *{Modifying with a constant value}

Replace the target of a \textbf{Lens} or all of the targets of a
\textbf{Setter} or \textbf{Traversal} with a constant value.

\begin{minted}{haskell}
(.~) :: ASetter s t a b -> b -> s -> t

> [1,2,3] & element 0 .~ 3
[3,2,3]
> 0 & bitAt 8 .~ True
256
> [1,2,3] & traversed . filtered odd .~ 0
[0,2,0]
\end{minted}
\end{box1}

\begin{box2}
\subsection *{Prisms and Isos}

An \textbf{Iso} is a pair of inverse functions. You can invert an \textbf{Iso} with \textbf{from}.

\textbf{Prism}s can be thought of as \textbf{Iso}s that can fail in one
direction. You can invert a \textbf{Prism} with \textbf{re}.

\begin{minted}{haskell}
type Prism s t a b 
  forall p f. (Choice p, Applicative f) =>
              p a (f b) -> p s (f t)
type Prism' s a = Prism s s a a

prism  :: (b -> t)
       -> (s -> Either t a)
       -> Prism s t a b
prism' :: (a -> s)
       -> (s -> Maybe a)
       -> Prism' s a

> 5^.re _Left ^?! _Left
5
> _Left # 1
Left 1

type Iso s t a b =
  forall p f. (Profunctor p, Functor f) =>
              p a (f b) -> p s (f t)
type Iso' s a = Iso s s a a

iso :: (s -> a) -> (b -> t) -> Iso s t a b
from :: AnIso s t a b -> Iso b a t s

> 'a' ^. from enum
97
> 97 ^. enum :: Char
'a'

> Map.empty & at "hi"
>           . non Map.empty
>           . at "world" ?~ "!"
fromList [("hi",fromList [("world","!")])]
\end{minted}

\end{box2}

\end{multicols}

\begin{box2}
\subsection *{Some setting operators}
\begin{tabular}{ l l l l l }
\textbf{Operator}  &  \textbf{W/result}  &  \textbf{W/state}  &  \textbf{W/both}  &  \textbf{Action}\\
\hline
\verb!+~ !         &  \verb!<+~ !        &  \verb!+= !        &  \verb!<+= !      &  Add to target(s)\\
\verb!-~ !         &  \verb!<-~ !        &  \verb!-= !        &  \verb!<-= !      &  Subtract from target(s)\\
\verb!*~ !         &  \verb!<*~ !        &  \verb!*= !        &  \verb!<*= !      &  Multiply target(s)\\
\verb!//~!         &  \verb!<//~!        &  \verb!//=!        &  \verb!<//=!      &  Divide target(s)\\
\verb!^~ !         &  \verb!<^~ !        &  \verb!^= !        &  \verb!<^= !      &  Raise target(s) to a non-negative Integral power\\
\verb!^^~!         &  \verb!<^^~!        &  \verb!^^=!        &  \verb!<^^=!      &  Raise target(s) to an Integral power\\
\verb!**~!         &  \verb!<**~!        &  \verb!**=!        &  \verb!<**=!      &  Raise target(s) to an arbitrary power\\
\verb!||~!         &  \verb!<||~!        &  \verb!||=!        &  \verb!<||=!      &  Logically or target(s)\\
\verb!&&~!         &  \verb!<&&~!        &  \verb!&&=!        &  \verb!<&&=!      &  Logically and target(s)\\
\verb!<>~!         &  \verb!<<>~!        &  \verb!<>=!        &  \verb!<<>=!      &  mappend to the target monoidal value(s)\\
\end{tabular}
\end{box2}


\end{document}
